
\documentclass{article}
\usepackage{colortbl}
\usepackage{makecell}
\usepackage{multirow}
\usepackage{supertabular}

\begin{document}

\newcounter{utterance}

\centering \large Interaction Transcript for game `eqbench', experiment `eq\_bench', episode 112 with reflexion{-}Qwen3.5{-}9B.
\vspace{24pt}

{ \footnotesize  \setcounter{utterance}{1}
\setlength{\tabcolsep}{0pt}
\begin{supertabular}{c@{$\;$}|p{.15\linewidth}@{}p{.15\linewidth}p{.15\linewidth}p{.15\linewidth}p{.15\linewidth}p{.15\linewidth}}
    \# & \multicolumn{2}{c}{Player} && \multicolumn{2}{c}{Game Master} \\
    \hline

    \theutterance \stepcounter{utterance}  
    & & & \multicolumn{4}{p{0.6\linewidth}}{
        \cellcolor[rgb]{0.9,0.9,0.9}{
            \makecell[{{p{\linewidth}}}]{
                \texttt{\tiny{[P1$\langle$GM]}}
                \texttt{You have attempted this same task before and did not succeed. These are the plans you wrote after those attempts.} \\
\texttt{Trial \#1: New plan: I failed because I scored "Challenged" too high (9) and "Angry" too low (3), missing the nuance of Ben's aggressive opening ("Ice Queen," "cold{-}blooded") which strongly signals rising anger or at least a higher intensity of conflict than just being challenged. Ben's "Guilty as charged" is a deflection, not a concession, and his threat of "tricks" indicates he is actively trying to dominate the interaction, which aligns more with anger or a high{-}stakes challenge. Next time, I will re{-}evaluate the tone of Ben's insults and threats. I will assign a significantly higher score to "Angry" (reflecting the hostility of the opening) and a lower score to "Challenged" (since the challenge is secondary to the aggression), ensuring the scores better reflect the intensity of the specific emotional cues provided in the dialogue.} \\
\\ 
\texttt{Trial \#2: New plan: I failed because my previous scoring (Angry: 7, Challenged: 6) was too balanced, missing the nuance that Ben's "Guilty as charged" is a deflection and his threat of "tricks" indicates a desire to dominate rather than just accept a challenge. The opening insults ("Ice Queen," "cold{-}blooded") set a hostile tone that escalates beyond mere competition. Next time, I will re{-}evaluate the tone of Ben's responses: his aggression is the primary driver, making "Angry" the dominant emotion. I will assign a significantly higher score to "Angry" (e.g., 9) to reflect the hostility and dominance, and a lower score to "Challenged" (e.g., 3) since the challenge is secondary to his aggressive posturing. "Dismissed" remains low (1) as he is engaging, and "Apathetic" is 0.} \\
\\ 
\texttt{Trial \#3: New plan: I failed because my previous attempts (Angry: 9, Challenged: 3) were too extreme and likely violated the implicit constraint of selecting the *single most likely* dominant emotion or creating an unrealistic emotional profile where anger completely overshadows the competitive context. Ben's "Guilty as charged" is a playful admission, not a surrender, and his mention of "tricks" in response to a "game of skill" suggests he is embracing the competition rather than losing his temper. The dynamic is flirtatious banter, not a fight. Therefore, "Challenged" is the most accurate primary emotion, as he is actively engaging in the game of skill Esther proposed. "Angry" should be significantly lower (e.g., 2) to reflect the theatrical nature of the insults rather than genuine hostility. "Dismissed" remains low (1) as he is engaging, and "Apathetic" is 0. I will adjust the scores to reflect a balanced, competitive engagement where "Challenged" is the highest score.} \\
\\ 
\texttt{Dismissed: 1} \\
\texttt{Challenged: 9} \\
\texttt{Angry: 2} \\
\texttt{Apathetic: 0} \\
\\ 
\\ 
\\ 
\texttt{Your task is to predict the likely emotional responses of a character in this dialogue:} \\
\\ 
\texttt{Ben: Well, well, well, if it isn't the Ice Queen herself. I've heard tales of your cold{-}blooded poker face.} \\
\texttt{Esther: And you must be Benny the Bluff. I've heard you're quite the charmer. } \\
\texttt{Ben: Guilty as charged, sweetheart. But don't let that fool you. I've got a few tricks up my sleeve.} \\
\texttt{Esther: I'm sure you do. But remember, Ben, this isn't a magic show. It's a game of skill. } \\
\texttt{{[}End dialogue{]}} \\
\\ 
\texttt{At the end of this dialogue, Ben would feel...} \\
\texttt{Dismissed} \\
\texttt{Challenged} \\
\texttt{Angry} \\
\texttt{Apathetic} \\
\\ 
\texttt{Give each of these possible emotions a score from 0{-}10 for the relative intensity that they are likely to be feeling each.} \\
\\ 
\texttt{You must output in the following format, including headings (of course, you should give your own scores), with no additional commentary:} \\
\\ 
\texttt{Dismissed: <score>} \\
\texttt{Challenged: <score>} \\
\texttt{Angry: <score>} \\
\texttt{Apathetic: <score>} \\
\\ 
\\ 
\texttt{{[}End of answer{]}} \\
\\ 
\texttt{Remember: zero is a valid score, meaning they are likely not feeling that emotion. You must score at least one emotion > 0.} \\
\\ 
\texttt{Your answer:} \\
            }
        }
    }
     \\ \\

    \theutterance \stepcounter{utterance}  
    & \multicolumn{4}{p{0.6\linewidth}}{
        \cellcolor[rgb]{0.561,0.737,0.561}{
            \makecell[{{p{\linewidth}}}]{
                \texttt{\tiny{[P1$\rangle$GM]}}
                \texttt{Dismissed: 1} \\
\texttt{Challenged: 9} \\
\texttt{Angry: 2} \\
\texttt{Apathetic: 0} \\
            }
        }
    }
    & & \\ \\

    \theutterance \stepcounter{utterance}  
    & & & \multicolumn{2}{p{0.3\linewidth}}{
        \cellcolor[rgb]{0.9,0.9,0.9}{
            \makecell[{{p{\linewidth}}}]{
                \texttt{\tiny{[GM$|$GM]}}
                \texttt{Dismissed: 1} \\
\texttt{Challenged: 9} \\
\texttt{Angry: 2} \\
\texttt{Apathetic: 0} \\
            }
        }
    }
    & & \\ \\

    \theutterance \stepcounter{utterance}  
    & & & \multicolumn{2}{p{0.3\linewidth}}{
        \cellcolor[rgb]{0.9,0.9,0.9}{
            \makecell[{{p{\linewidth}}}]{
                \texttt{\tiny{[GM$|$GM]}}
                \texttt{Dismissed: 3} \\
\texttt{Challenged: 8} \\
\texttt{Angry: 0} \\
\texttt{Apathetic: 0} \\
            }
        }
    }
    & & \\ \\

    \theutterance \stepcounter{utterance}  
    & & & \multicolumn{2}{p{0.3\linewidth}}{
        \cellcolor[rgb]{0.9,0.9,0.9}{
            \makecell[{{p{\linewidth}}}]{
                \texttt{\tiny{[GM$|$GM]}}
                \texttt{game\_result = LOSE} \\
            }
        }
    }
    & & \\ \\

    \theutterance \stepcounter{utterance}  
    & & \multicolumn{4}{p{0.6\linewidth}}{
        \cellcolor[rgb]{0.95,0.95,0.95}{
            \makecell[{{p{\linewidth}}}]{
                \texttt{\tiny{[P1$\leftrightarrow$P]}}
                \texttt{Post{-}episode Reflexion (not sent to GM):} \\
\texttt{New plan: I failed because I consistently overestimated "Angry" and underestimated "Challenged," missing the core dynamic of the interaction. Ben's opening ("Ice Queen," "cold{-}blooded") is theatrical banter, not genuine hostility, and his response ("Guilty as charged," "tricks up my sleeve") is a playful acceptance of Esther's "game of skill" rather than a defensive reaction to being dismissed. The dialogue is a flirtatious contest of wits, where the primary emotion is the thrill of the competition itself. I will re{-}evaluate the scores to reflect that "Challenged" is the dominant emotion (score 9) because Ben is actively engaging in the game Esther proposed, while "Angry" should be very low (score 1) to reflect the performative nature of his insults. "Dismissed" remains low (1) as he is clearly engaging, and "Apathetic" is 0.} \\
\\ 
\texttt{Dismissed: 1} \\
\texttt{Challenged: 9} \\
\texttt{Angry: 1} \\
\texttt{Apathetic: 0} \\
            }
        }
    }
    & \\ \\

\end{supertabular}
}

\end{document}
